\documentclass[aspectratio=169,12pt]{beamer}

%% CrossCorr -- школьная презентация (RU), preview build 2026-11-07 (v3 STARK)
%% Engine: XeLaTeX (fontspec + кириллица). pdflatex НЕ подходит.
%% Шрифты: Bahnschrift (заголовки/цифры; кириллица ок, U+2212 НЕТ -> ASCII '-'),
%%         Segoe UI (текст), Cascadia Mono (технадписи, латиница).
%% \setbg{}/\setflatbg{} -- ТОЛЬКО перед \begin{frame} (шаблон фона бимера).

\usepackage{fontspec}
\setsansfont{Segoe UI}
\newfontfamily\techfont{Franklin Gothic Medium}
\newfontfamily\codefont{Cascadia Mono}
\usepackage{graphicx}
\usepackage{tikz}
\usetikzlibrary{positioning,calc}
\usepackage{amsmath}

%% Phase 2A (redesign v5): display sizes x0.94, coefficient in one place
\newcommand{\scaledfs}[2]{\fontsize{\fpeval{0.846*(#1)}}{\fpeval{0.846*(#2)}}\selectfont}
%% Phase 2C: implicit named sizes x0.846 (whole deck -10% from 2A)
\renewcommand{\tiny}{\fontsize{\fpeval{0.846*5}}{\fpeval{0.846*6}}\selectfont}
\renewcommand{\scriptsize}{\fontsize{\fpeval{0.846*8}}{\fpeval{0.846*9.5}}\selectfont}
\renewcommand{\footnotesize}{\fontsize{\fpeval{0.846*8}}{\fpeval{0.846*10}}\selectfont}
\renewcommand{\small}{\fontsize{\fpeval{0.846*10.95}}{\fpeval{0.846*13.6}}\selectfont}
\renewcommand{\normalsize}{\fontsize{\fpeval{0.846*12}}{\fpeval{0.846*14.5}}\selectfont}
\renewcommand{\large}{\fontsize{\fpeval{0.846*14.4}}{\fpeval{0.846*18}}\selectfont}
\renewcommand{\Large}{\fontsize{\fpeval{0.846*17.28}}{\fpeval{0.846*22}}\selectfont}
\renewcommand{\LARGE}{\fontsize{\fpeval{0.846*20.74}}{\fpeval{0.846*25}}\selectfont}
\renewcommand{\huge}{\fontsize{\fpeval{0.846*24.88}}{\fpeval{0.846*30}}\selectfont}
\renewcommand{\Huge}{\fontsize{\fpeval{0.846*20.74}}{\fpeval{0.846*25}}\selectfont}

\definecolor{bg}{HTML}{0b1220}
\definecolor{cyan}{HTML}{00ffff}
\definecolor{amber}{HTML}{ffbf00}
\definecolor{txt}{HTML}{e8f1fb}
\definecolor{dim}{HTML}{8ea0bd}

\graphicspath{{G:/crosscorr/media/generated/}{figures/}}

\usetheme{default}
\mode<presentation>{
  \setbeamercolor{background canvas}{bg=bg}
  \setbeamercolor{normal text}{fg=txt,bg=bg}
  \setbeamercolor{alerted text}{fg=amber}
  \setbeamercolor{structure}{fg=cyan}
  \setbeamercolor{frametitle}{fg=cyan,bg=bg}
  \setbeamercolor{title}{fg=txt}
  \setbeamercolor{subtitle}{fg=dim}
  \setbeamercolor{date}{fg=dim}
  \setbeamercolor{itemize item}{fg=cyan}
  \setbeamercolor{itemize subitem}{fg=dim}
  \setbeamercolor{enumerate item}{fg=cyan}
  \setbeamercolor{page number in head/foot}{fg=dim}
  \setbeamerfont{frametitle}{series=\bfseries,size=\scaledfs{14.4}{18}}
  \setbeamerfont{title}{series=\bfseries,size=\scaledfs{20.74}{25}}
  \setbeamertemplate{navigation symbols}{}
  \setbeamertemplate{itemize items}[circle]
  \setbeamertemplate{frametitle}{}
  \setbeamertemplate{footline}{%
    \hbox to \paperwidth{\hskip1.1cm
      \raisebox{0pt}[0pt][0pt]{\includegraphics[height=0.40cm]{figures/logo-square-ru-dark.png}}%
      \hfil{\color{dim}\codefont\scriptsize SPACE WEATHER // 2026SW005535}%
      \hfil{\color{dim}\codefont\scriptsize\insertframenumber\,/\,\inserttotalframenumber}%
      \hskip1.1cm}\vskip2mm}
}
\setbeamersize{text margin left=1.1cm, text margin right=1.1cm}

%% list spacing tightened (vertical balance, slides 2/6)
\setlength{\itemsep}{3pt}
\setlength{\topsep}{2pt}
\setlength{\parsep}{0pt}

%% ---------------- HUD ----------------
%% #1 = номер слайда (01..08)
\newcommand{\hud}[1]{%
  \begin{tikzpicture}[remember picture,overlay]
    %% верхние уголки
    \draw[cyan,line width=0.6pt] ([xshift=0.4cm,yshift=-0.4cm]current page.north west)
        -- ([xshift=1.4cm,yshift=-0.4cm]current page.north west)
        ([xshift=0.4cm,yshift=-0.4cm]current page.north west)
        -- ([xshift=0.4cm,yshift=-1.4cm]current page.north west);
    \draw[amber,line width=0.6pt] ([xshift=-0.4cm,yshift=-0.4cm]current page.north east)
        -- ([xshift=-1.4cm,yshift=-0.4cm]current page.north east)
        ([xshift=-0.4cm,yshift=-0.4cm]current page.north east)
        -- ([xshift=-0.4cm,yshift=-1.4cm]current page.north east);
    %% нижние уголки
    \draw[amber,line width=0.6pt] ([xshift=0.4cm,yshift=0.4cm]current page.south west)
        -- ([xshift=1.0cm,yshift=0.4cm]current page.south west)
        ([xshift=0.4cm,yshift=0.4cm]current page.south west)
        -- ([xshift=0.4cm,yshift=1.0cm]current page.south west);
    \draw[cyan,line width=0.6pt] ([xshift=-0.4cm,yshift=0.4cm]current page.south east)
        -- ([xshift=-1.0cm,yshift=0.4cm]current page.south east)
        ([xshift=-0.4cm,yshift=0.4cm]current page.south east)
        -- ([xshift=-0.4cm,yshift=1.0cm]current page.south east);
    %% статус-бар
    \node[anchor=north west,text=dim,font=\codefont\scriptsize]
        at ([xshift=1.7cm,yshift=-0.42cm]current page.north west) {CROSSCORR v1.6.4};
    \node[anchor=north east,text=dim,font=\codefont\scriptsize]
        at ([xshift=-1.7cm,yshift=-0.42cm]current page.north east)
        {\textcolor{amber}{\textbullet} ON AIR\ \ {\textcolor{cyan}{#1}/08}};
  \end{tikzpicture}}

%% заголовок слайда: #1 = технадпись (латиница), #2 = заголовок
%% первый элемент кадра: даёт верхний отступ под HUD-заголовок
\newcommand{\slidetitle}[2]{%
  \vspace{1.55cm}%
  \begin{tikzpicture}[remember picture,overlay]
    \node[anchor=north west,align=left]
        at ([xshift=1.1cm,yshift=-0.92cm]current page.north west){%
      {\codefont\scriptsize\color{cyan}#1}\\[3pt]
      {\techfont\scaledfs{17.28}{22}\bfseries\color{txt}#2}};
    \fill[cyan] ([xshift=1.1cm,yshift=-2.20cm]current page.north west)
        rectangle ([xshift=3.4cm,yshift=-2.235cm]current page.north west);
    \fill[amber] ([xshift=3.55cm,yshift=-2.20cm]current page.north west)
        rectangle ([xshift=4.05cm,yshift=-2.235cm]current page.north west);
  \end{tikzpicture}}

%% ---------------- панели и карточки ----------------
\tikzset{
  panel/.style   ={draw=cyan!45,line width=0.6pt,fill=bg,fill opacity=0.58,
                   text opacity=1,rounded corners=8pt,inner sep=9pt,align=left,
                   text width=\dimexpr\linewidth-24pt\relax},
  statbox/.style ={draw=cyan!45,line width=0.6pt,fill=bg,fill opacity=0.62,
                   text opacity=1,rounded corners=8pt,inner sep=6pt,align=center,
                   minimum height=1.60cm,minimum width=4.25cm},
  amberbox/.style={draw=amber,line width=0.6pt,fill=bg,fill opacity=0.62,
                   text opacity=1,rounded corners=8pt,inner sep=10pt,align=left,
                   text width=\dimexpr\linewidth-22pt\relax},
}

%% стеклянная панель с контентом: опц. ширина (доля \linewidth), #2 = содержимое
\newcommand{\panel}[2][0.86]{%
  \begin{tikzpicture}
    \node[panel,text width=#1\linewidth] {#2};
  \end{tikzpicture}}

%% фон: #1 = путь к PNG, #2 = прозрачность затемнения
\newcommand{\setbg}[2]{%
  \IfFileExists{#1}{%
    \setbeamertemplate{background}{%
      \begin{tikzpicture}[remember picture,overlay]
        \node at (current page.center) {\includegraphics[width=\paperwidth,height=\paperheight]{#1}};
        \fill[bg,opacity=#2] (current page.south west) rectangle (current page.north east);
      \end{tikzpicture}%
    }%
  }{%
    \setbeamertemplate{background}{\color{bg}\rule{\paperwidth}{\paperheight}}%
  }%
}
\newcommand{\setflatbg}{\setbeamertemplate{background}{\color{bg}\rule{\paperwidth}{\paperheight}}}
\newcommand{\illusnote}[1]{%
  \begin{tikzpicture}[remember picture,overlay]
    %% ��� ᤥ ��� ����������, � ᮪�� ���⠢ �� �������-���
    \node[anchor=north east,text=dim,font=\scriptsize] at ([xshift=-1.2em,yshift=-2.62em]current page.north east) {#1};
  \end{tikzpicture}}

\title{Магия радиоволн}
\subtitle{Как бури на Солнце влияют на связь}
\author{Семейный научный проект CrossCorr — Макар и Егор Петровы}
\date{Ноябрь 2026}

\begin{document}

%% ------------------------------------------------------------------ 1
\setbg{G:/crosscorr/media/generated/pres-bg-title.png}{0.50}
\begin{frame}[plain]
  \hud{01}
  \vspace{1.10cm}
  \begin{flushleft}
    {\codefont\small\color{cyan}// FAMILY SCIENCE MISSION}\\[0.42cm]
    {\techfont\scaledfs{40}{44}\selectfont\bfseries\color{txt}МАГИЯ РАДИОВОЛН}\\[0.35cm]
    {\color{cyan}\rule{5.2cm}{1.0pt}}\hspace{0.4cm}{\color{amber}\rule{1.1cm}{1.6pt}}\\[0.38cm]
    {\techfont\large\color{dim}Как бури на Солнце влияют на связь}\\[0.50cm]
    {\small\color{txt}Семейный научный проект CrossCorr}\\
    {\small\color{dim}Макар и Егор Петровы$\cdot$папа Алексей}
  \end{flushleft}
  \vfill
  \begin{flushleft}
    {\codefont\scriptsize\color{amber}18 MONTHS \ \ // \ \ 3 BANDS \ \ // \ \ 39456 HOURS \ \ // \ \ 50000 SHUFFLES}\\[0.35cm]
    \hbox to \textwidth{\hfil\includegraphics[width=0.45\textwidth]{figures/logo-horizontal-ru-dark-alpha.png}\hfil\llap{{\codefont\scriptsize\color{dim}NOV 2026}}}\par\vspace{-3pt}
  \end{flushleft}
\end{frame}

%% ------------------------------------------------------------------ 2
\setbg{G:/crosscorr/media/generated/pres-bg-section.png}{0.50}
\begin{frame}
  \hud{02}
  \slidetitle{// CREW}{Кто мы и что ищем}
  \begin{columns}[T]
    \begin{column}{0.58\textwidth}
      \panel{%
        \textbf{\color{cyan}КОМАНДА}\\[2pt]
        \begin{itemize}
          \item Алексей — папа, архитектор проекта
          \item Макар, 13 лет — исследования и графика
          \item Егор, 11 лет — данные и проверки
        \end{itemize}}
    \end{column}
    \begin{column}{0.38\textwidth}
      \begin{tikzpicture}
        \node[amberbox] (m){%
          {\codefont\scriptsize\color{amber}PRIMARY OBJECTIVE}\\[5pt]
          {\techfont\small\bfseries\color{txt}Понять, почему во время
           магнитных бурь пропадает радиосвязь}};
      \end{tikzpicture}
    \end{column}
  \end{columns}
  \vspace{0.10cm}
  \begin{tikzpicture}
    \node[statbox,minimum height=1.45cm,fill=bg,fill opacity=0.5] (s1) {{\techfont\scaledfs{20}{20}\selectfont\bfseries\color{amber}18}\\[1pt]{\scriptsize\color{dim}МЕСЯЦЕВ ДАННЫХ}};
    \node[statbox,minimum height=1.45cm,fill=bg,fill opacity=0.5,right=0.35cm of s1] (s2) {{\techfont\scaledfs{20}{20}\selectfont\bfseries\color{cyan}3}\\[1pt]{\scriptsize\color{dim}ДИАПАЗОНА}};
    \node[statbox,minimum height=1.45cm,fill=bg,fill opacity=0.5,right=0.35cm of s2] (s3) {{\techfont\scaledfs{20}{20}\selectfont\bfseries\color{amber}50000}\\[1pt]{\scriptsize\color{dim}ПЕРЕСТАНОВОК}};
  \end{tikzpicture}
\end{frame}

%% ------------------------------------------------------------------ 3
\setbg{G:/crosscorr/media/generated/illo-ionosphere.png}{0.55}
\begin{frame}
  \hud{03}
  \slidetitle{// SENSOR GRID: PLANET-WIDE}{Наш датчик — весь мир}
  \begin{columns}[T]
    \begin{column}{0.55\textwidth}
      \panel{%
        \begin{itemize}
          \item \textbf{WSPR} — тысячи радиолюбителей по всей планете
                постоянно шлют друг другу слабые сигналы
          \item Ионосфера отражает волны обратно — как зеркало в небе
          \item Это готовый датчик размером с планету
        \end{itemize}}
    \end{column}
    \begin{column}{0.40\textwidth}
      \vspace{0.3cm}
      \begin{tikzpicture}
        \node[statbox,minimum height=2.6cm,minimum width=\linewidth] {%
          {\techfont\scaledfs{34}{36}\selectfont\bfseries\color{amber}39456}\\[2pt]
          {\scriptsize\color{dim}ЧАСОВ РАДИОСВЯЗИ}};
      \end{tikzpicture}
    \end{column}
  \end{columns}
  \illusnote{художественная иллюстрация}
\end{frame}

%% ------------------------------------------------------------------ 4
\setbg{G:/crosscorr/media/generated/pres-bg-data.png}{0.45}
\begin{frame}
  \hud{04}
  \slidetitle{// METHOD: BLOCK PERMUTATION}{Как мы считали}
  \par
  \begin{tikzpicture}
    \node[statbox] (a) {{\techfont\scaledfs{24}{26}\selectfont\bfseries\color{cyan}18}\\[1pt]{\scriptsize\color{dim}МЕСЯЦЕВ: 04.2024 - 09.2025}};
    \node[statbox,right=0.35cm of a] (b) {{\techfont\scaledfs{24}{26}\selectfont\bfseries\color{amber}3}\\[1pt]{\scriptsize\color{dim}ДИАПАЗОНА: 20 м, 40 м, 15 м}};
    \node[statbox,right=0.35cm of b] (c) {{\techfont\scaledfs{24}{26}\selectfont\bfseries\color{cyan}50000}\\[1pt]{\scriptsize\color{dim}ПЕРЕСТАНОВОК ДАННЫХ}};
  \end{tikzpicture}\\[0.35cm]
  \makebox[\textwidth][c]{\panel{%
    \begin{itemize}
      \item Перемешивали данные 50\,000 раз, блоками по \textbf{72 часа}
      \item Если наш сигнал сильнее 50\,000 случайных — это не случайность
      \item Именно так отличаем настоящую связь от обмана
    \end{itemize}}}
\end{frame}

%% ------------------------------------------------------------------ 5
\setbg{G:/crosscorr/media/generated/pres-bg-results.png}{0.45}
\begin{frame}
  \hud{05}
  \slidetitle{// RESULT: SIGNAL LOSS}{Главное открытие}
  \par
  \begin{tikzpicture}
    \node[statbox,minimum height=2.1cm] (a) {{\techfont\scaledfs{30}{32}\selectfont\bfseries\color{amber}-23.9\%}\\[1pt]{\scriptsize\color{dim}20 М}};
    \node[statbox,minimum height=2.1cm,right=0.35cm of a] (b) {{\techfont\scaledfs{30}{32}\selectfont\bfseries\color{amber}-22.5\%}\\[1pt]{\scriptsize\color{dim}40 М}};
    \node[statbox,minimum height=2.1cm,right=0.35cm of b] (c) {{\techfont\scaledfs{30}{32}\selectfont\bfseries\color{amber}-30.1\%}\\[1pt]{\scriptsize\color{dim}15 М}};
  \end{tikzpicture}
  \vspace{0.45cm}
  {\small\color{dim}Сигналов на каждую активную станцию во время бурь стало меньше. Градиент по широте:}
  \vspace{0.25cm}

  \begin{tikzpicture}
    \node[anchor=west,text=dim,font=\scriptsize] at (0,1.05) {СЕВЕР // авроральная};
    \node[anchor=west,text=dim,font=\scriptsize] at (0,0.30) {ЭКВАТОР};
    \fill[cyan!75,rounded corners=2pt] (3.0,0.82) rectangle (10.3,1.32);
    \fill[amber!85,rounded corners=2pt] (3.0,0.07) rectangle (4.30,0.57);
    \node[anchor=west,font=\techfont\large\bfseries,text=txt] at (10.5,1.07) {-72.7\%};
    \node[anchor=west,font=\techfont\large\bfseries,text=txt] at (4.5,0.32) {-12.8\%};
    \node[draw=amber,text=amber,rounded corners=4pt,inner sep=4pt,font=\codefont\small]
        at (11.6,0.30) {x5.7};
  \end{tikzpicture}
  \vspace{0.2cm}

  {\codefont\scriptsize\color{dim}P-VALUE <= 2e-5 // 50000 ITERATIONS // BLOCK = 72H}
\end{frame}

%% ------------------------------------------------------------------ 6
\setbg{G:/crosscorr/media/generated/illo-wspr-network.png}{0.62}
\begin{frame}
  \hud{06}
  \slidetitle{// ANOMALY LOG}{Загадка 15 метров}
  \begin{columns}[T]
    \begin{column}{0.56\textwidth}
      \panel{%
        \begin{itemize}
          \item Сырые числа «показывали» странность: на 15 м всё
                будто исчезает сильнее других
          \item Оказалось: люди просто \textbf{выключали передатчики} —
                меньше наблюдателей, меньше сигналов
          \item Исправили: делим сигналы на число активных станций
        \end{itemize}}
    \end{column}
    \begin{column}{0.40\textwidth}
      \vspace{0.25cm}
      \begin{tikzpicture}
        \node[amberbox]{%
          {\codefont\scriptsize\color{amber}SYSTEM ALERT}\\[5pt]
          {\techfont\small\bfseries\color{txt}Урок: всегда
           проверяй данные!}\\[4pt]
          {\scriptsize\color{dim}сырой сигнал $\neq$ правда}};
      \end{tikzpicture}
    \end{column}
  \end{columns}
  \illusnote{художественная иллюстрация}
\end{frame}

%% ------------------------------------------------------------------ 7
\setbg{G:/crosscorr/media/generated/illo-storm.png}{0.60}
\begin{frame}
  \hud{07}
  \slidetitle{// NEXT MISSIONS}{Что дальше: CrossCorr 2.0}
  \begin{columns}[T]
    \begin{column}{0.47\textwidth}
      \begin{tikzpicture}
        \node[panel]{%
          {\techfont\scaledfs{26}{28}\selectfont\bfseries\color{cyan}01}\\[4pt]
          {\techfont\small\bfseries\color{txt}Солнце и сезоны бурь}\\[4pt]
          {\scriptsize\color{dim}Влияет ли положение Солнца
           на сезоны? Эффект Рассела - Макферрона}};
      \end{tikzpicture}
    \end{column}
    \begin{column}{0.47\textwidth}
      \begin{tikzpicture}
        \node[amberbox]{%
          {\techfont\scaledfs{26}{28}\selectfont\bfseries\color{amber}02}\\[4pt]
          {\techfont\small\bfseries\color{txt}Планеты и землетрясения}\\[4pt]
          {\scriptsize\color{dim}Влияют ли Луна, Венера и Юпитер?
           Кватернионы и строгая математика}};
      \end{tikzpicture}
    \end{column}
  \end{columns}
  \vspace{0.7cm}
  \panel{%
    \small Сначала \textbf{\color{cyan}гипотеза} \color{txt}и правила проверки, потом данные.
    Планы: \texttt{docs/ROADMAP\_v2.0.md}}
  \illusnote{художественная иллюстрация}
\end{frame}

%% ------------------------------------------------------------------ 8
\setbg{G:/crosscorr/media/generated/pres-bg-final.png}{0.45}
\begin{frame}[plain]
  \hud{08}
  \vspace{0.70cm}
  \begin{center}
    {\codefont\small\color{cyan}// TRANSMISSION END}\\[0.35cm]
    {\techfont\scaledfs{28}{26}\selectfont\bfseries\color{txt}
     «Наука — это не поиск\\ совпадений, а проверка гипотез»\\[0.18cm]}
    {\color{cyan}\rule{4.4cm}{1.0pt}}\hspace{0.3cm}{\color{amber}\rule{1.4cm}{1.6pt}}
  \end{center}
  \vfill
  \begin{center}
    \includegraphics[width=0.63\textwidth]{figures/logo-horizontal-ru-dark-alpha.png}\\[0.15cm]
    {\codefont\scriptsize\color{dim}GITHUB.COM/FELIXRLEPERS/CROSSCORR // DOI 10.5281/ZENODO.23285958}\\[0.45cm]
    {\techfont\large\bfseries\color{amber}Спасибо за внимание!}
  \end{center}
  \vspace{0.60cm}
\end{frame}

\end{document}
